%------------------------
% phdtaketaketake — PhD Application CV Template
%
% Adapted from a working physics-PhD-applicant CV. Keep the structure;
% replace the <angle-bracket> placeholders with your information. The
% file as-shipped compiles to a placeholder "demo CV" so you can verify
% the pipeline works (LaTeX install + custom resume commands) BEFORE
% personalizing. After verifying the pipeline, edit this file:
%
%   phdtaketaketake-cv-template --print > cv.tex
%   phdtaketaketake-cv-compile cv.tex                 # produces cv.pdf
%   # ... edit cv.tex to replace placeholders ...
%   phdtaketaketake-cv-compile cv.tex                 # re-compile
%
% Section policy:
%
%   - All sections are present by default. Delete any you do not need
%     (entire \section{...} block + its \resumeSubHeadingListStart/End
%     wrapper). Empty section headers look worse than missing sections.
%   - Sections that should not be deleted (always relevant for PhD apps):
%     Education, Research Experience.
%   - Sections frequently deleted: Teaching, Leadership.
%   - Some content is wrapped in OPTIONAL_BLOCK markers — off by default,
%     uncomment line-by-line to enable. Currently:
%       * Relevant Coursework table (under Education)
%       * Multi-affiliation extra emails (under Header)
%
% OPTIONAL_BLOCK convention: the descriptive prose that explains *when*
% to enable a block lives ABOVE the OPTIONAL_BLOCK_START marker. The
% block itself contains only LaTeX code (each line prefixed with `% `),
% so "uncommenting all lines between START and END" gives you valid
% LaTeX without leaking description text into the document body.
%
% Tailoring for a target PI / program: see references/cv_optimization.md
% §"Tailoring playbook". The skill never invents content — tailoring is
% reordering and user-approved pruning only.
%------------------------

\documentclass[a4paper,11pt]{article}

% --- Required packages ---
% Most are in any standard TeX install (BasicTeX, MacTeX, TeX Live,
% MiKTeX). The three exceptions, marked below, ship with the
% "texlive-latex-extra" / "texlive-fonts-extra" bundles on Debian-style
% distros — minimal-TeX users (e.g. TinyTeX) may need to install them
% explicitly with `tlmgr install titlesec enumitem fancyhdr`.
\usepackage[a4paper,top=0.5in,bottom=0.5in,left=0.5in,right=0.5in]{geometry}
\usepackage{titlesec}      % (texlive-latex-extra) — \titleformat for section styling
\usepackage[usenames,dvipsnames]{color}
\usepackage{enumitem}      % (texlive-latex-extra) — leftmargin / label opts on itemize
\usepackage[pdftex]{hyperref}
\usepackage{fancyhdr}      % (texlive-latex-extra) — \pagestyle{fancy} + clearing headers/footers

\pagestyle{fancy}
\fancyhf{}
\fancyfoot{}
\renewcommand{\headrulewidth}{0pt}
\renewcommand{\footrulewidth}{0pt}

\urlstyle{rm}

\raggedbottom
\raggedright
\setlength{\tabcolsep}{0in}

% Sections formatting
\titleformat{\section}{
  \vspace{-10pt}\scshape\raggedright\large
}{}{0em}{}[\color{black}\titlerule \vspace{-6pt}]

%-------------------------
% Custom resume commands — DO NOT EDIT BELOW THIS LINE UNTIL \begin{document}

\newcommand{\resumeItemWithoutTitle}[1]{
  \item\small{
    #1 \vspace{-2pt}
  }
}

\newcommand{\resumeSubheading}[4]{
  \vspace{-1pt}\item
    \begin{tabular*}{0.97\textwidth}{l@{\extracolsep{\fill}}r}
      \textbf{#1} & \footnotesize #2 \\
      \footnotesize \textit{#3} & \footnotesize \textit{#4} \\
    \end{tabular*}\vspace{-5pt}
}

\renewcommand{\labelitemii}{\textbullet}

\newcommand{\resumeSubHeadingListStart}{\begin{itemize}[leftmargin=0.15in, label={}]}
\newcommand{\resumeSubHeadingListEnd}{\end{itemize}}
\newcommand{\resumeItemListStart}{\begin{itemize}}
\newcommand{\resumeItemListEnd}{\end{itemize}\vspace{-5pt}}

%-----------------------------
%%%%%%  CV STARTS HERE  %%%%%%

\begin{document}

%----------HEADING-----------------
%
% Optional content: extra emails for users with multiple affiliations
% (REU host institution, summer-program domain, dual-affiliated lab).
% Add as many lines as you need; each renders as a separate email
% under the primary one. To enable: uncomment the lines between
% OPTIONAL_BLOCK_START and OPTIONAL_BLOCK_END below.
%
\begin{tabular*}{\textwidth}{l@{\extracolsep{\fill}}r}
  \textbf{\LARGE <Your Full Name>} &
  \begin{tabular}[t]{r}
    Email: \href{mailto:<primary@example.edu>}{<primary@example.edu>} \\
% --- OPTIONAL_BLOCK_START: Multi-affiliation extra emails ---
%    \href{mailto:<secondary@example.edu>}{<secondary@example.edu>} \\
%    \href{mailto:<tertiary@example.edu>}{<tertiary@example.edu>} \\
% --- OPTIONAL_BLOCK_END: Multi-affiliation extra emails ---
    Website: \url{<https://your-site.example.com>}
  \end{tabular}
\end{tabular*}

%-----------EDUCATION-----------------
\section{Education}
  \resumeSubHeadingListStart
    \resumeSubheading
      {<University Name>}{<Start Date> - <End Date or "Expected June YYYY">}
      {<Degree, e.g., Bachelor of Science in Physics>; GPA: <X.XX>}{<City, State>}

% Optional content: Relevant Coursework table.
% Useful when applying to programs that explicitly weigh coursework
% (most physics / applied math / theory CS PhDs). Less useful for
% experimental bio / chem / wet-lab disciplines, where coursework is
% assumed. To enable: uncomment the lines between OPTIONAL_BLOCK_START
% and OPTIONAL_BLOCK_END below; the leading `\item[]` is required so
% the tabular attaches to a fresh (label-less) itemize entry.
%
% --- OPTIONAL_BLOCK_START: Relevant Coursework ---
%    \item[]
%    \vspace{4.5mm}
%    \noindent Relevant Coursework: \\
%    \vspace{1.5em}
%    \begin{tabular}{@{\hspace{1.5em}} p{0.45\linewidth} p{0.45\linewidth}}
%        \textbullet\ <Course 1> & \textbullet\ <Course 4> \\
%        \textbullet\ <Course 2> & \textbullet\ <Course 5> \\
%        \textbullet\ <Course 3> & \textbullet\ <Course 6> \\
%    \end{tabular}
% --- OPTIONAL_BLOCK_END: Relevant Coursework ---

  \resumeSubHeadingListEnd

% Add a second institution (study-abroad, transfer, master's) here by
% copying the \resumeSubheading{...}{...}{...}{...} pattern above into a
% new \resumeSubHeadingListStart / End wrapper.

\vspace{-5pt}
\section{Research Experience}
  \resumeSubHeadingListStart

    \resumeSubheading
    {<Research Project Title>}{<Start Date> - <End Date or "Present">}
    {Mentor: <Prof. Name>, <Institution>}{}
    \resumeItemListStart
        \resumeItemWithoutTitle
          {<First bullet — what you built / proved / measured. Lead with a verb. Quantify when possible (sample size, sensitivity, fold improvement).>}
        \resumeItemWithoutTitle
          {<Second bullet — methods, tools, scale. Name specific software / techniques (e.g., Geant4, MadGraph5, MCMC, Tensor Hadronic Decay analysis).>}
        \resumeItemWithoutTitle
          {<Third bullet — outcome / contribution. Cite paper / talk / preprint if applicable.>}
    \resumeItemListEnd
    \vspace{0.5em}

    % --- Add more research experiences below by copying the
    %     \resumeSubheading + \resumeItemListStart/End pattern above. ---
    %
    % Order: most relevant to your target PI / program FIRST. When
    % tailoring against a phdtaketaketake match.json, lead with the
    % experience whose topic / methods / system overlap most with the
    % top-bucket PIs' `research_areas`. See
    % references/cv_optimization.md §"Tailoring playbook".

\resumeSubHeadingListEnd

\vspace{0.1em}
\section{Publications \& Presentations}
% Single-layer itemize (Publications has no \resumeSubheading — items
% are contributed directly via \item[] sub-headings + \resumeItemWithoutTitle).
\resumeSubHeadingListStart

    \item[]\textbf{SELECTED PAPERS}
    \vspace{2pt}

    \resumeItemWithoutTitle
    {<Author list with \textbf{Your Name} bolded>. ``<Paper title>.'' \href{<https://doi-or-arxiv-url>}{\textit{<Venue, vol., year>}} (<year>).}

    % --- Add more selected papers above by copying the
    %     \resumeItemWithoutTitle{...} line. Delete the entire
    %     SELECTED PAPERS sub-block (item[] header + items) if you
    %     have no published / accepted papers yet. ---

    \vspace{1em}

    \item[] \textbf{WORK IN PROGRESS}
    \vspace{2pt}

    \resumeItemWithoutTitle
    {<Author list with \textbf{Your Name}>. ``<Working title>.'' \textit{In Preparation} (<expected year>).}

    % --- Add more work-in-progress papers above by copying the line.
    %     Delete the entire sub-block if you have no in-prep work. ---

    \vspace{1em}

    \item[] \textbf{POSTERS \& TALKS}
    \vspace{2pt}

    \resumeItemWithoutTitle
    {\textbf{<Your Name>}. ``<Poster / talk title>.'' \textit{<Symposium / conference name>} (<Month Year>).}

    % --- Add more posters / talks above by copying the line.
    %     Delete the entire sub-block if not applicable. ---

\resumeSubHeadingListEnd
\vspace{-5pt}

\section{Technical Skills}
% Single-layer itemize: items go directly inside \resumeSubHeadingListStart.
\resumeSubHeadingListStart
    \item\small{\textbf{<Skill Category 1, e.g. Programming \& Machine Learning>}: <comma-separated list, e.g. Python, PyTorch, NumPy, SciPy>}
    \vspace{-2pt}
    \item\small{\textbf{<Skill Category 2, e.g. Physics Software>}: <comma-separated list, e.g. Geant4, MadGraph5, Pythia8, Delphes>}
    \vspace{-2pt}
    \item\small{\textbf{Languages}: <e.g. English (Professional Proficiency), Mandarin (Native)>}
\resumeSubHeadingListEnd

\vspace{-5pt}
\section{Teaching Experience}
  \resumeSubHeadingListStart
    \resumeSubheading
    {<Role, e.g., Learning Assistant: Classical Mechanics>}{<Start Date> - <End Date or "Present">}
    {<Institution>}{}
    \resumeItemListStart
        \resumeItemWithoutTitle{<One bullet — what you did and the scale (number of students / sessions / weeks). Lead with a verb.>}
    \resumeItemListEnd
\resumeSubHeadingListEnd

% Delete the entire Teaching Experience section above (this comment +
% \section{...} + the \resumeSubHeadingListStart/End wrapper) if you
% have no teaching experience yet.

\vspace{-5pt}
\section{Leadership Experience}
  \resumeSubHeadingListStart
    \resumeSubheading
    {\href{<https://org-website.example.com>}{<Organization Name>}}{<Start Date> - <End Date>}
    {<Role, e.g., Webmaster / President / Treasurer>}{<City, State>}
    \resumeItemListStart
        \resumeItemWithoutTitle
        {<One bullet — impact and scope of the role. Avoid passive phrasing ("was responsible for"); lead with a verb.>}
    \resumeItemListEnd
\resumeSubHeadingListEnd

% Delete the entire Leadership Experience section above if you have no
% leadership / org / club roles yet.

\vspace{-5pt}
\section{Honors and Awards}
% Single-layer itemize: items go directly inside \resumeSubHeadingListStart.
% No leading \vspace before the first \item — that triggers a "missing \item" error.
\resumeSubHeadingListStart
    \item\small{<Award Name (Year) — \$amount or selectivity note (e.g. "1 of 12 nationally") if relevant>}
    \vspace{-2pt}
    \item\small{<Award Name (Year)>}
    \vspace{-2pt}
    \item\small{<Award Name (Year)>}
\resumeSubHeadingListEnd

\end{document}
